\subsection{All complex divisor coordinates and exact norm compatibility}

\begin{statement}{E.72}{Proved}{Surviving prime chains, principal-class equations, and saturation of the positional norm system}
Keep the actual fields and roots
\[
 \begin{gathered}
 K=\mathbb Q(\zeta),\qquad
 L_d=K(\beta_d),\quad \beta_d^{3^d}=17\zeta,\qquad
 F_d=\mathbb Q(\alpha_d),\quad \alpha_d^{3^d}=17,\\
 \beta_{d+1}^3=\beta_d,\qquad\alpha_{d+1}^3=\alpha_d .
 \end{gathered}
 \tag{E72.1}
\]
Thus $L_0=K$ and $F_0=\mathbb Q$. For a number field $E$,
write $\widehat{E^\times}_3=\varprojlim_r E^\times/(E^\times)^{3^r}$
and
\[
 \widehat{\operatorname{Div}(E)}_3
  =\varprojlim_r\bigoplus_{\mathfrak p}\mathbb Z/3^r,\qquad
 \mathscr R(J)=\{(a_j)\in\prod_{j\in J}\mathbb Z_3:
       a_j\longrightarrow0\pmod{3^r}\text{ for every }r\}.
 \tag{E72.2}
\]
These divisor groups carry their displayed prodiscrete topology.
In particular, the infinite restricted sequence factor is not
asserted compact.

Here is the entire complex divisor norm limit. For each rational
prime $p\ne3,17$ such that
$v_3(\operatorname{ord}_p(17))=1$, there is exactly one prime
$\mathfrak p\mid p$ of $K$ for which the residue of $17\zeta$
has order prime to three. Call the set of these primes $\mathcal P_1$.
Each $p$ is split in $K$. For $\mathfrak p\in\mathcal P_1$, let
$\mathfrak P_{d,\mathfrak p}$ be the degree-one prime of $L_d$
specified by
\[
 \beta_d\equiv r_{d,\mathfrak p}\pmod{\mathfrak P_{d,\mathfrak p}},
 \qquad
 r_{d,\mathfrak p}^{\,3^d}=\overline{17\zeta},
 \qquad r_{d,\mathfrak p}\text{ has order prime to three}.
 \tag{E72.3}
\]
The root in (E72.3) is unique in the prime-to-three subgroup of
$\mathbb F_p^\times$; its cube is the corresponding root at the
preceding level. Let $\mathfrak q_{3,d}$ and $\mathfrak q_{17,d}$
be the unique primes of $L_d$ above $3$ and $17$. The explicit map
\[
 \begin{split}
 \mathcal D&=\mathbb Z_3^2\oplus\mathscr R(\mathcal P_1)
       \ \xrightarrow{\ \sim\ }\
       \varprojlim_{d,N}\widehat{\operatorname{Div}(L_d)}_3,\\
 a=(a_3,a_{17},(a_{\mathfrak p}))
       &\longmapsto
 D_d(a)=a_3\mathfrak q_{3,d}
       +a_{17}\mathfrak q_{17,d}
       +\sum_{\mathfrak p\in\mathcal P_1}
                      a_{\mathfrak p}\mathfrak P_{d,\mathfrak p}
 \end{split}
 \tag{E72.4}
\]
is a topological isomorphism. Its inverse reads the coefficients
of the divisor at $K$. No other prime chain contributes a
nonzero compatible completed divisor. The further rational norm
sends these coordinates to
\[
 a_3[3]+2a_{17}[17]
       +\sum_{\mathfrak p\in\mathcal P_1}
                              a_{\mathfrak p}[p].
 \tag{E72.5}
\]
Consequently this rational divisor map is injective.

For every $d$, let $C_d=\operatorname{Cl}(L_d)\{3\}$, and use
the actual classes of the displayed ideals to define
\[
 \begin{split}
 \mathcal D_{\mathrm{pr}}
   &=\left\{a\in\mathcal D:
       a_3[\mathfrak q_{3,d}]
       +\sum_{\mathfrak p\in\mathcal P_1}
                  a_{\mathfrak p}[\mathfrak P_{d,\mathfrak p}]=0
                       \text{ in }C_d\text{ for every }d\right\},\\
 \mathcal X&=\varprojlim_{d,N}\widehat{L_d^\times}_3,\qquad
 \mathcal E_\infty=
       \varprojlim_{d,N}
                   \bigl(\mathcal O_{L_d}^{\times}\otimes\mathbb Z_3\bigr).
 \end{split}
 \tag{E72.6}
\]
The term at $17$ is absent from the class equation because
$(\beta_d)=\mathfrak q_{17,d}$. The restricted sums in the finite
groups $C_d$ are well defined. These are literal principal-class
equations in the fields
$X^{2\cdot3^d}+17X^{3^d}+289$, with the ideals specified in
(E72.3)--(E72.4).

The entire completed multiplicative norm system fits into the
exact sequence
\[
 0\longrightarrow\mathcal E_\infty
  \longrightarrow\mathcal X
  \xrightarrow{\ \operatorname{div}\ }\mathcal D_{\mathrm{pr}}
  \longrightarrow0,\qquad
 \mathcal E_\infty\simeq
                  \prod_{d\ge0}\mathbb Z_3^{\,2\cdot3^d}.
 \tag{E72.7}
\]
In particular the class equations in (E72.6) are sufficient for
one coherent multiplicative lift, including compatibility at all
levels. Every nonempty fiber of the divisor map is a compact
torsor under the displayed product. Its product coordinates are
the actual recursive norm-kernel coordinates of E.69.

The base element, including its root-of-unity factor, has a
finite-level test. For $a\in\mathcal D_{\mathrm{pr}}$, choose a
completed principal generator $g_1$ with
$\operatorname{div}(g_1)=D_1(a)$, using E.57's specified ideal
coordinates, and put
\[
                    \nu(a)=N_{L_1/K}(g_1).
 \tag{E72.8}
\]
This is independent of the generator, because the completed
ordinary-unit norm from $L_1$ is $1$. It is an injective map with
divisor $D_0(a)$, and the base projection of (E72.7) is exactly
$\nu\circ\operatorname{div}$. Thus, for an actual
$z\in\widehat{K^\times}_3$, membership in
$\operatorname{pr}_0(\mathcal X)$ is equivalent to the following
three explicit requirements: its divisor has the support in
(E72.4); that divisor satisfies every class equation (E72.6);
and $z=\nu(\operatorname{div}z)$. The last equality tests the
remaining $\mu_3$ factor by E.57's norm formulas. The higher
class matrices and the integer $h$ of E.69 remain unevaluated
arithmetic data; this statement is an exact membership theorem,
not a numerical evaluation of those matrices.

The prime above three gives a particularly precise interpretation
of that remaining integer. If
$N_{L_d/K}\mathcal U_d=\lambda^{3^{a_d}\mathbb Z_3}$
as in E.69, then
\[
 \operatorname{ord}[\mathfrak q_{3,d}]=3^{a_d},\qquad
 (\beta_1+3\zeta+1)\mathcal O_{L_d}
                =\mathfrak q_{3,d}^{\,3^{d-1}}
                    \quad(d\ge1).
 \tag{E72.9}
\]
The order here is the order in the full ideal class group;
it is a power of three. In particular
$[\mathfrak q_{3,2}]$ has exact order three. The stable integer
$h$ is the eventual exponent of these specified prime classes.
There is no assertion that equality at two adjacent levels
proves stabilization.

There are also exact saturation and positional comparison maps.
Every member of $\mathcal X$ whose rational base norm is supported
on $\{3,17\}$ consists of $\{3,17\}$-units at every level.
The entire such subgroup is
\[
 \begin{split}
 \mathcal X_{\{3,17\}}
   &=\{(\beta_d^{\,b}u_d)_d:
               b\in\mathbb Z_3,\ (u_d)\in\mathcal U_\infty\},\\
 \operatorname{pr}_0(\beta^b u)
   &=(17\zeta)^b\lambda^{3^h z},\qquad
 N_{K/\mathbb Q}\operatorname{pr}_0(\beta^b u)
                  =17^{2b}3^{3^h z}.
 \end{split}
 \tag{E72.10}
\]
The inverse removes the common valuation $b$ at the unique prime
above $17$. In particular rational norm supported only at $3$
forces $b=0$ and gives exactly the whole family of E.69.
For positional tuples with pure component $1$, the total-norm
cube condition is $b\in3\mathbb Z_3$. Writing $b=3r$, their
scaled toric coordinates from E.46 are
\[
       \left(\kappa_3(17),\
             2r\,\kappa_3(17)+3^{h-1}z\,\kappa_3(3)\right).
 \tag{E72.11}
\]

For the full positional algebra
$R_{d+1}=F_d\times L_d$, put
$\mathcal P_0=\{p\ne3,17:v_3(\operatorname{ord}_p17)=0\}$.
The pure divisor limit away $3,17$ has exactly one degree-one
prime chain for each $p\in\mathcal P_0$. Its coefficient norm
is the identity. The sets $\mathcal P_0$ and the rational
primes underlying $\mathcal P_1$ are disjoint. At $3$ the
pure chain has two persistent branches, both of residue
degree one; the complex chain has one. At $17$ the pure
and complex chains have residue degrees one and two.
In these actual branch coordinates the rational divisor norm is
\[
 \begin{array}{c|c}
 \text{prime}&\text{norm on the compatible coefficients}\\ \hline
 3&(b_3,c_3,a_3)\longmapsto b_3+c_3+a_3\\
 17&(e_{17},f_{17})\longmapsto e_{17}+2f_{17}\\
 p\in\mathcal P_0\cup\mathcal P_1&a_p\longmapsto a_p .
 \end{array}
 \tag{E72.12}
\]
Its kernel is exactly $\mathbb Z_3^3$, supported at $3$ and $17$.
Therefore every full positional norm tuple whose scaled second
toric coordinate is supported on $\{3,17\}$ consists of actual
$\{3,17\}$-unit components in both field towers. This is the
map to the quotient $H$ of E.46; it makes no assertion of a
lift through the intrinsic metabelian or marked comparison kernels.
The manuscript anchor is \S51.16, fols.\ 733v--736v,
scans 172--178, through the actual covering and toric maps of
E.25, E.27 and E.46.
\end{statement}

\begin{proof}
First fix a prime $\mathfrak p\nmid3\cdot17$ of $K$, with residue
size $q$. The roots $\beta_d\mu_{3^d}$ identify the root set with
$\mathbb Z/3^d$. The actual arithmetic Frobenius acts by
\[
              x\longmapsto qx+B\pmod{3^d},
 \tag{E72.13}
\]
where the compatible $B\in\mathbb Z_3$ is the local Kummer
cocycle of $17\zeta$ at the chosen Frobenius. A different
root system changes $B$ by $(q-1)t$, transporting all the
following prime chains by the same translation.
Frobenius orbits are precisely the primes above $\mathfrak p$,
and their lengths are their relative residue degrees. If a
higher orbit lies over a lower orbit, the ideal norm multiplies
its coefficient by the ratio of those lengths.

The affine transformation in (E72.13) has an integral fixed
point exactly when $-B/(q-1)\in\mathbb Z_3$. This is equivalent
to the residue of $17\zeta$ belonging to the prime-to-three
subgroup of $\mathbb F_q^\times$: the local principal-unit
group has bijective cubing, and the finite residue group
supplies precisely the remaining Kummer class. When it
exists, the fixed point is unique, since the integer $q>1$
is not $1$ in $\mathbb Z_3$.

Every other orbit has unbounded three-part length. For a
direct verification, write the fixed point in $\mathbb Q_3$
as $x_*=-B/(q-1)$. Iteration gives
$\phi^j(x)-x=(q^j-1)(x-x_*)$. Since $q\equiv1\pmod3$,
$v_3(q^{3^k}-1)=v_3(q-1)+k$. On a fixed clopen residue
class which does not contain an integral fixed point,
$v_3(x-x_*)$ is bounded. Thus the orbit lengths of every
refinement have three-parts tending uniformly to infinity.
The norm coefficient of a compatible divisor on that lower
orbit is divisible by every power of three and is zero.
This reasoning uses a finite sum at each higher level,
so it also applies to completed divisors.

Only the singleton fixed-point branch can remain. Its residue
degrees are all one and its norm coefficient is constant.
Conversely any coefficient in $\mathbb Z_3$ defines a compatible
divisor on that branch. This proves the asserted local inverse
limit, together with its literal coefficient map.

For an inert rational $p\ne3,17$, the element $17$ lies in
$\mathbb F_p^\times$, of order prime to three, whereas
$\zeta$ has order three in $\mathbb F_{p^2}^\times$.
There is no complex fixed branch. For split $p$, the two
residues of the radicand are $17\zeta_{\mathfrak p}$ and
$17\zeta_{\mathfrak p}^{-1}$. Their ratio has order three,
so at most one has order prime to three. Such a residue
exists precisely when the three-primary component of $17$
has exact order three; its inverse determines the required
value of $\zeta_{\mathfrak p}$. This proves both the stated
set $\mathcal P_1$ and formula (E72.3).

At $3$ and $17$ the complex tower has the unique primes and
total ramification established in E.64; their relative residue
degrees are one. The completed divisor coefficients at all
other primes have just been determined. At each level the
condition that they tend to zero modulo every $3^r$ is exactly
the same condition on the surviving base coefficients.
This proves (E72.4) as an isomorphism of the specified
prodiscrete groups. The residue degrees from $K$ to $\mathbb Q$
are respectively $1,2,1$, giving (E72.5) and its injectivity.

For clarity, the completion and the inverse limit in (E72.6)
are handled separately. For any fixed number field $E$, the
discrete divisor sequence gives
\[
 0\longrightarrow
    \mathcal O_E^\times\otimes\mathbb Z_3
 \longrightarrow \widehat{E^\times}_3
 \xrightarrow{\ \operatorname{div}\ }
    \widehat{\operatorname{Div}(E)}_3
 \longrightarrow \operatorname{Cl}(E)\{3\}
 \longrightarrow0.
 \tag{E72.14}
\]
Indeed the group of principal divisors is a subgroup of the
free divisor group and is itself free. The sequence defining
it therefore splits as a sequence of abstract abelian groups
before completion. Its induced topology inside the divisor
group is the multiplication topology: if the three-primary
class group has exponent $3^e$ and a principal divisor is
$3^{r+e}D$, the prime-to-three part of $[D]$ is already zero,
and $3^e[D]=0$. Hence the intersection with
$3^{r+e}\operatorname{Div}(E)$ is contained in $3^r$ times
the principal-divisor group. This proves
the exact image and the stated topology in (E72.14).
The kernel is compact because the ordinary-unit group is
finitely generated.

Necessity of every class equation is immediate from (E72.14).
For sufficiency, fix $a$ satisfying all of them. At each $d$
let $T_d$ be the fiber of the divisor map in (E72.14) over
$D_d(a)$. It is a nonempty compact torsor under the ordinary-unit
completion. Norms send $T_{d+1}$ to $T_d$.
In the compact product $\prod_dT_d$, impose the closed
equations $N(x_{d+1})=x_d$. Every finite collection of them
has a solution: choose an element in the highest required
fiber, norm it down to all preceding fibers, and choose
arbitrary elements in the other fibers. The finite intersection
property gives a simultaneous solution. This proves the
surjectivity in (E72.7), without assuming surjectivity of
any individual ordinary-unit norm or compactness of the
whole multiplicative completion.

The kernel is precisely $\mathcal E_\infty$. An ordinary-unit
norm tuple is a $3$-unit tuple in E.69. Its base lies in
$\mu_3$ and also in
$I_0=\lambda^{3^h\mathbb Z_3}$, so it is $1$.
Conversely, in a coherent $3$-unit tuple with base $1$,
the coefficient at the unique prime above $3$ is constant
under norm and therefore zero at every level. Such a tuple
consists of ordinary units. Consequently
$\mathcal E_\infty=\ker(\mathcal U_\infty\to I_0)$,
and E.69 gives exactly the product in (E72.7) with its
recursive coordinates. Translation by any one lift proves
the asserted compact topology and coordinates on each fiber.

Any two choices of $g_1$ in (E72.8) differ by a completed
ordinary unit. E.57 computes its norm as $1$, including
the torsion generator. Thus $\nu$ is well defined.
Every coherent lift of $a$ has first component in this same
fiber, so its base is $\nu(a)$. Conversely every $a$ in
$\mathcal D_{\mathrm{pr}}$ has such a lift, already proved.
Its divisor at $K$ is exactly the base coordinate $a$.
This proves injectivity, the entire base-image criterion,
and the claimed finite test for the root factor. E.57's
principal generators and norm matrix specify the calculation;
no assertion that the divisor alone determines an arbitrary
element of $\widehat{K^\times}_3$ has been used.

Let $m_d$ be the full order of $[\mathfrak q_{3,d}]$.
The included element $c_1=\beta_1+3\zeta+1$ has ideal
$\mathfrak q_{3,d}^{3^{d-1}}$, since its ideal at $L_1$
is $\mathfrak q_{3,1}$ and every adjacent extension is
totally ramified. Hence $m_d$ is a power of three.
The image of the valuation
$\mathcal O_{L_d}[1/3]^\times\to\mathbb Z$
is exactly $m_d\mathbb Z$: this is the definition of the
order of that ideal class. Its completed image is
$m_d\mathbb Z_3$. The relative residue degree at $3$ is one,
and E.69 places every total norm in
$\lambda^{\mathbb Z_3}$. Reading its $\lambda$-valuation
therefore gives $m_d=3^{a_d}$ and (E72.9).
E.69 supplies $a_2=1$ and eventual equality $a_d=h$.

If the rational base norm of a coherent complex tuple has
no valuation outside $\{3,17\}$, the injectivity and explicit
support in (E72.5) force every other compatible divisor
coefficient to vanish. Exactness in (E72.14), or its version
with these two primes inverted, then places each component
in the completed $\{3,17\}$-unit group. Its common
coefficient $b$ at $17$ is removed by the literal coherent
tuple $(\beta_d^b)_d$, leaving exactly a member of
$\mathcal U_\infty$. This proves (E72.10).
The rational norm is as displayed because
$N_{K/\mathbb Q}\lambda=3$ and
$N_{K/\mathbb Q}(17\zeta)=17^2$.
The cube condition and the division by three in E.46 then
give (E72.11).

For the pure comparison use the actual root set
$\alpha_d\mu_{3^d}$ over $\mathbb Q$ at $p\ne3,17$.
Its affine Frobenius is $x\mapsto px+B$.
If $p\equiv1\pmod3$, the argument above applies unchanged,
with radicand $17$ and fixed branch exactly when
$v_3(\operatorname{ord}_p17)=0$.
If $p\equiv2\pmod3$, $p-1$ is a unit in $\mathbb Z_3$,
so there is always a unique integral fixed point.
For the other orbits apply the same length argument to
Frobenius squared: $p^2\equiv1\pmod3$ and
$v_3(p^{2\cdot3^k}-1)=v_3(p^2-1)+k$.
This proves the pure assertion and disjointness of the
surviving prime sets.

Finally the local factors at $3$ are the actual factors
of E.65. If $\rho^3=17$ is the root in $\mathbb Q_3$ and
$m=3^{d-1}$, then for $d\ge1$
\[
 X^{3^d}-17=(X^m-\rho)(X^{2m}+\rho X^m+\rho^2).
 \tag{E72.15}
\]
Both factors become Eisenstein at $3$ after $X\mapsto X-1$.
Cubing carries each factor to its predecessor, and each
successive extension on either branch is totally ramified
of degree three. Both residue degrees are one, including
their maps to $\mathbb Q_3$. At $17$, the pure polynomial
is Eisenstein with residue degree one, whereas the complex
branch has the residue degree two already used above.
This proves (E72.12). A basis of its kernel is
$(1,-1,0)$ and $(1,0,-1)$ at $3$, and $(-2,1)$ at $17$.
There is no kernel away these primes. The total-norm map
of E.46 is precisely this rational norm on divisor
coordinates, so the final saturation assertion follows.

There is a concrete finite obstruction to prolonging the
first norm correction in E.57. At $\pi=2-\zeta$ above $7$,
the roots of $\beta_1^3=17\zeta$ are $3,5,6$ in
$\mathbb F_7$. The only root of order prime to three is
$6=-1$. The two other primes are inert cubic in $L_2/L_1$,
since $3$ and $5$ are not cubes in $\mathbb F_7^\times$.
E.57's element $w=(g_1/\pi)^2$ has the actual divisor
\[
 \operatorname{div}(w)=
 4\mathfrak p_{\beta=3}
       -2\mathfrak p_{\beta=6}
       -2\mathfrak p_{\beta=5}.
 \tag{E72.16}
\]
The coefficients at the two inert primes are not divisible
by three. Consequently $w$ is not a norm from
$\widehat{L_2^\times}_3$, even though its first norm is
the root $\zeta$. The obstruction is the actual ideal-norm
map at $7$; it requires no inference from a nonprincipal
ideal or an uncomputed higher class matrix.

The retained finite certificate compares the actual
factorizations of both defining polynomials with the
affine-Frobenius orbit formulas. It also checks the weighted
norm transition matrices, the eligible-prime sets, and
the displayed saturated kernel. These finite checks support
the formulas; the uniform orbit-length proof above supplies
the all-level assertion.
\end{proof}
